in degree @@M0@@.

The module of differentials @@M1@@ is isomorphic to @@M2@@, since one can define an @@M3@@-derivation @@M4@@:
@@M5@@
so as to satisfy the universal property defining @@M6@@, by setting @@M7@@; moreover, @@M8@@ is compatible with the gradings if the factor @@M10@@ is assigned degree @@M9@@.

Let @@M11@@ be the vector bundle over @@M12@@ defined by @@M13@@. Then @@M14@@ is the restriction to @@M15@@ of @@M16@@. Now @@M17@@ is the affine @@M18@@-scheme defined by @@M19@@, so @@M20@@ is defined by @@M21@@ (since @@M22@@ is locally free of finite type), hence @@M23@@ is equal to @@M24@@, and consequently @@M25@@. To see its grading, one looks locally, and sees at once that @@M26@@ must be regarded as being of degree @@M27@@.

Taking the degree @@M28@@ terms in the exact sequence (5.2.1) gives the announced exact sequence.

\sgaStableTarget{sga6.e08.rem.5-2-1}\textbf{Remark 5.2.1.} Using the general results on the differential study of Hilbert functors ([2], 5), one directly obtains an isomorphism @@M29@@ from the given functorial characterization of @@M30@@ for every @@M31@@-morphism @@M32@@.

\sgaStableTarget{sga6.e08.para.5-3}\textbf{5.3.} By 5.2, we have:
@@M33@@
Consequently:
@@M34@@
